\section*{\inTextCircled[10]{5}{15} Doors to Aberdour}\stepcounter{section}\phantomsection\addcontentsline{toc}{section}{Doors to Aberdour}
{\entryfont The ancient transport road continues until it reaches a broad junction carved into the surrounding rock. The main tunnel carries on straight ahead, while a second route branches sharply to the right toward a monumental gateway.}
\begin{DndReadAloud}
	The passage widens into a great junction.

	Straight ahead, the old transport road continues into darkness, its course eventually broken by the jagged remains of a collapsed bridge spanning a vast chasm.

	To your right, another route rises toward an enormous metal gate set into the surrounding stone. Its surface is covered in weathered geometric engravings and ancient Dwarven script.
\end{DndReadAloud}
\subsection*{\inTextCircled[10]{5A}{10} Dwarven Door}\stepcounter{subsection}\phantomsection\addcontentsline{toc}{subsection}{Dwarven Door}
{\entryfont The route ends before a massive gate forged from a single expanse of dark, ancient metal. Rows of Dwarven runes cover its surface, though many have faded almost beyond recognition.

A squat stone pedestal stands several feet before the door. Several differently shaped metal inlets surround a shallow central bowl, suggesting that the mechanism was intended to be operated from here.}
\begin{DndReadAloud}
	The passage terminates before an immense metal door set deep into the surrounding rock. Its dark surface is covered in rows of weathered Dwarven runes, some still sharply engraved, others reduced to little more than faded scars.

	Before it stands a broad stone pedestal. A shallow bowl occupies its centre, surrounded by several metal-lined recesses of different colours and finishes.

	Whatever mechanism once opened the gate appears to remain intact.	
\end{DndReadAloud}
\subsubsection*{Opening the Door}
{\entryfont The door is sealed by an ancient Dwarven locking mechanism operated through the pedestal. Opening it requires the characters to correctly interact with the objects and inlets arranged upon it.

The exact puzzle can be replaced with another suitable challenge if desired. One possible implementation, \hyperref[ex:DwarfPuzzle]{\LinkFont{A Dwarf Knows His Craft}}, is described in the Resources Chapter at the end of this document.}
\subsection*{\inTextCircled[10]{5B}{10} Meeting Ralathor}\stepcounter{subsection}\phantomsection\addcontentsline{toc}{subsection}{Meeting Ralathor}
{\entryfont Beyond the Dwarven gate, the passage continues only a short distance before ending at a second sealed entrance. Unlike the first, this door is covered in far more intricate mechanisms, delicate fittings, and unfamiliar geometric patterns.}
\begin{DndReadAloud}
	The passage beyond the first gate is quiet, disturbed only by the faint clicking of old machinery and the low murmur of a voice ahead.

	A solitary man stands before the second door, leaning close to its surface as he studies one of the many interlocking mechanisms set into it.

	"No... that cannot be right."

	He adjusts something, waits, and receives no response.

	Only then does he notice your approach.

	He straightens slowly and looks between you, the open Dwarven gate behind you, and the constructs standing before him.

	"Interesting", he says. "You are certainly not who I expected to find down here."
\end{DndReadAloud}
{\noindent\entryfont Ralathor introduces himself only as a mysterious hermit travelling through the ancient tunnels. He is cautious but curious, and quickly begins asking the expedition about the world above.

His knowledge of the tunnels is incomplete. Ralathor knows that the ancient Dwarves constructed an extensive subterranean network beneath much of the Kingdom of Fife, and he suspects that similar passages may extend throughout other parts of Caledonia. He does not know where every route leads or what settlements may once have existed along them.

The door before him is particularly puzzling. Its mechanisms are far more intricate than the Dwarven engineering he expected to encounter within the tunnels, and its construction appears fundamentally different from the gate the party has just opened. Ralathor cannot explain who built it or why.}
\subsubsection*{A Sudden Change in Priorities}
{\entryfont Ralathor initially questions the expedition out of curiosity: where they came from, what brought them into the tunnels, and what is happening in the world above.

Once he learns of Zargothrax, the fall of Dundee, and the coming war, his demeanour changes immediately. His attention shifts away from the door and toward the route behind the party. He begins asking much more practical questions:
\begin{itemize}
	\item "You came from Auchtertool?"
	\item "How far?"
	\item "And Prince Angus - where is he now?"
\end{itemize}
Ralathor quickly becomes preoccupied with determining the fastest route back to the surface and how he might reach Prince Angus McFife. He gives no clear explanation for his urgency, only insisting that he needs to speak with the prince as soon as possible.}
\subsection*{\inTextCircled[10]{5C}{10} Gnomish Door}\stepcounter{subsection}\phantomsection\addcontentsline{toc}{subsection}{Gnomish Door}
{\entryfont The second gate bears little resemblance to the monumental Dwarven door behind the expedition. Its surface is covered in fine metalwork, recessed panels, rotating fittings, and several small mechanisms whose purposes are not immediately apparent.}
\begin{DndReadAloud}
	The passage ends before another sealed door, but this one is unlike anything you have seen in the tunnels so far.

	Brass fittings, narrow grooves, small dials, and interlocking mechanisms cover much of its surface. Several components appear movable, while others vanish into the stonework around the frame.

	Nothing about it suggests a simple lock.
\end{DndReadAloud}
{\noindent\entryfont An example puzzle used to open the door is described in Resources: \hyperref[ex:GnomePuzzle]{\LinkFont{What Makes it Tick?}}}
\subsubsection*{Ralathor's Help}
{\entryfont Ralathor has already spent some time examining the door before the expedition arrives. Although he has not discovered its solution, he has identified several properties of the mechanism and can provide hints if the players become stuck.

These hints should reveal how individual parts of the puzzle behave rather than provide the final solution. Once the door opens, Ralathor leaves the expedition behind and heads back toward Auchtertool, intent on reaching Prince Angus McFife as quickly as possible.}
\subsubsection*{Hidden Button}
{\entryfont A character who thoroughly examines the door or surrounding stonework may discover a small concealed button among the mechanisms.

Pressing it causes a series of engraved numbers around the button to illuminate, beginning at 20 and counting steadily downward. Pressing the button again resets the count to 20. Where a zero would normally appear is instead a small skull-and-crossbones symbol.

If the countdown reaches the skull, the final symbol illuminates, a mechanism gives a quiet click-

and nothing happens.}